\section{Introduction}

Weak intensity in a spin spectrum need not mean an absence of excitations.
In a system with several magnetic sites per unit cell, destructive interference
can suppress the matrix element of the physical probe while the collective
mode remains present. Cluster interference and symmetry-constrained
momentum dependences are established tools in magnetic
scattering~\cite{waldmann2003interference,haraldsen2011pentamer,tabrizi2021selection,tabrizi2021anisotropic}.
Effective chains of cluster spins likewise provide a familiar route to
low-energy modes and folded spectra, including trimer and tetramer
chains~\cite{cheng2022trimer,li2026tetramer}.

The projected form factor is only the leading approximation to a scattering
amplitude. Perturbative corrections to dimer structure factors already show
that interactions can modify this momentum dependence~\cite{mueller2003structure}.
Waldmann and G\"udel constructed effective scattering operators beyond simple
spin projection and analyzed tensor corrections caused by anisotropy-induced
mixing of cluster multiplets~\cite{waldmann2005manyspin}.
More recently, systematic extinctions have been distinguished from zeros that
depend on Hamiltonian parameters~\cite{tabrizi2026extinctions}.
Neither a form-factor zero nor the use of a corrected effective observable is
therefore the question addressed here.

The failure of symmetry to fix an intensity is also established.
For isotropic clusters, Ghassemi Tabrizi derives a universal momentum
dependence when a unique transition-operator irrep mediates a transition
and occurs once in the site-spin operator representation.
When the irrep occurs repeatedly, Eq.~(16) of
Ref.~\onlinecite{tabrizi2021selection} instead retains independent reduced
matrix elements and their interference. Symmetry therefore need not fix
their ratio. This does not decide whether the ratios for all thermally
accessible transitions can be adjusted together by physical exchange
controls. That simultaneous cancellation is the issue addressed below.

Temperature makes the interpretation of a dark point more demanding.
Effective-observable vertex corrections already describe conditional
excitations and thermal changes of spectral weight in spin
ladders~\cite{exius2010vertex}. Thermally activated multiparticle
transitions and spectral-weight transfer have also been analyzed in
frustrated ladders~\cite{honecker2016thermal}. Thermal lifting of an
accidental intensity cancellation has also been observed for interfering
nuclear and magnetovibrational phonon amplitudes~\cite{raymond2018interference}.
The appearance of additional lines on heating is therefore not itself a
new mechanism. The question
here is whether an exchange parameter can be retuned to cancel all those
transitions together, rather than only the ground-state response.

The Heisenberg pentamer chain provides doublet and quartet local ground
multiplets without changing the microscopic spin length or bond topology.
We show that their retuned zeros have different thermal stability. A
doublet permits cancellation of the entire first-order probe, whereas a
quartet produces a composite vector whose bond-spin channels cannot all be
canceled by the same real retuning. A local commutator argument establishes
this obstruction at nonzero scaled temperature for every finite even
periodic chain in the specified effective model. An exact second-moment
sum rule extends the algebraic obstruction to other local spins,
conditional on the same dipole--composite probe structure. Combined with
locality, it gives a size-independent lower bound on finite-frequency
weight in the uniform-trace limit. Allowing independent first-order
SU(2)-invariant retunings within each cell does not lower the optimal
response below uniform retuning: additional operator momenta contribute
nonnegative spectral weight. A four-cell ring is dark
at zero temperature but acquires an exactly calculable activated response.
Microscopic spectra test the weak-coupling powers and the four-cell activated
line groups. Finite-coupling backgrounds and finite momentum
acceptance limit the use of this distinction
as an experimental diagnostic.

\section{Model and low-energy observable}

Each cell contains five spin-$1/2$ sites at unit spacing. With $J_1=1$,
$\alpha=J_2/J_1$, and $\beta=J_3/J_1>0$, the Hamiltonian is
\begin{align}
 H_0&=\sum_{n=0}^{N-1}\big[
 \alpha(\bm S_{n,1}\!\cdot\!\bm S_{n,2}
       +\bm S_{n,4}\!\cdot\!\bm S_{n,5})\notag\\
 &\hspace{29mm}+\bm S_{n,2}\!\cdot\!\bm S_{n,3}
       +\bm S_{n,3}\!\cdot\!\bm S_{n,4}\big],\\
 H&=H_0+\beta V,\qquad
 V=\sum_{\langle n,n+1\rangle}\bm S_{n,5}\!\cdot\!\bm S_{n+1,1}.
 \label{eq:rank-ham}
\end{align}
The bond sum is specified as open or periodic in each calculation, and the
microscopic length is $L=5N$. We assume that the isolated cell has a single
ground multiplet of spin $s$, separated from all other multiplets by
$\Dloc>0$. The numerical examples below have $s=1/2$ or $3/2$.

Rotational symmetry within this irrep fixes the local projection,
\begin{equation}
 P_n\bm S_{n,a}P_n=c_a\bm T_n,\qquad
 c_a=c_{6-a},\qquad \sum_a c_a=1.
 \label{eq:rank-projection}
\end{equation}
Here $\bm T_n$ is a spin-$s$ operator. The leading Hamiltonian is an
antiferromagnetic effective chain with exchange $\beta c_1^2$.
This standard reduction does not by itself determine a suppressed intensity.

It is convenient to center the intra-cell Fourier phase:
\begin{align}
 B_n^z(q)&=\sum_{a=1}^5 e^{iq(a-3)}S_{n,a}^z,\\
 O_q&=\frac{1}{\sqrt{5N}}\sum_n e^{iQn}B_n^z(q),
 \quad Q=5q\pmod{2\pi}.
\end{align}
Changing the coordinate origin adds only a global phase to $O_q$ and does not
change a spectral weight. Equation~\eqref{eq:rank-projection} gives
\begin{align}
 P O_qP&=A(q,\alpha)D_Q,\notag\\
 A(q,\alpha)&=\sum_a c_a\cos[q(a-3)],\\
 D_Q&=\frac{1}{\sqrt{5N}}\sum_n e^{iQn}T_n^z.
 \label{eq:rank-amplitude}
\end{align}
The normalization of $D_Q$ is inherited from the microscopic probe, not
chosen independently for each curve.

Our observable is the integrated inelastic low-band weight
\begin{equation}
 \Wlow=\frac{1}{n_G}\sum_{r=1}^{n_G}
 \sum_{0<E_m-E_G<\Dloc/2}|\langle m|O_q|G_r\rangle|^2.
 \label{eq:rank-weight}
\end{equation}
All transitions within the ground eigenspace are excluded.
The gap $\Dloc$ in this cutoff and in $g=\beta/\Dloc$ is evaluated at the
unshifted local root. For each tested finite system we check that the cutoff
contains the complete retained low band. This is a fixed-size,
weak-coupling construction, not a claim of a uniformly isolated many-body
band as $N$ tends to infinity.
